\documentclass[preprint,12pt]{elsarticle}
\usepackage{amsmath,amssymb}
\usepackage{booktabs}
\usepackage{tabularx}
\usepackage{multirow}
\usepackage[hidelinks]{hyperref}
\setlength{\emergencystretch}{3em}
\journal{JCDL (primary target) / Data Analysis and Knowledge Discovery (fallback)}

\begin{document}

\begin{frontmatter}
\title{The Lexicon Is the Load-Bearing Wall: How Schema Term Quality Decides
Entity-Layer Fate in Dictionary-Driven Enrichment at Production Scale}
\author[addr1]{Zhiqian Chen\corref{cor1}}
\ead{chenzhiqian@cdut.edu.cn}
\cortext[cor1]{Corresponding author.}
\address[addr1]{Chengdu University of Technology, Chengdu, China}

\begin{abstract}
Dictionary-driven mention extraction is a common design in production
knowledge-base pipelines: it is cheap, reproducible, and idempotent. We
report a production incident in which this design, fed by a
machine-generated schema, silently collapsed the entity layer of a
61{,}137-document scholarly knowledge base into \textbf{575 entities}
carrying \textbf{19.2 million entity links}---every entity a bare type word
such as ``method'' or ``energy''. No conventional monitor fired: enrichment
jobs completed, throughput was normal, and link volume looked like growth.
We trace the collapse through three layers---the extraction mechanism (pure
lexicon scan), the lexicon's provenance (a batch-generated schema whose
per-type aliases form a category taxonomy rather than a domain terminology),
and the healthy counter-example in the same deployment (78 documents, a
hand-curated 3{,}742-term lexicon, 1{,}608 well-formed entities). We then
describe a three-source, anti-hallucination lexicon generation method
(transplanted seed terms, corpus-candidate classification, and per-type
expansion, all gated by corpus attestation; 3{,}702 and 1{,}279 terms for
the two affected schemas) and a full-chain rebuild over 766{,}501
re-enqueued jobs that doubled as a load test, surfacing and fixing five
pipeline pathologies (a whole-corpus join in graph projection that
saturated a 112-core database at load 500; two link-insertion deadlock
classes; an $O(N^2)$ acronym check amplified by the richer lexicon
itself---137$\times$ link throughput after the fixes; and a parameter-order
bug whose rows were silently swallowed by \texttt{INSERT IGNORE} for seven
hours while every job-level monitor read green). The rebuilt layer holds
\textbf{1{,}871 canonical entities with 25.1M links and 44k typed
relations} on the 61{,}137-document library---3.25$\times$ the collapsed
cardinality, a relation layer where none functioned, and scoped retrieval
coverage up +6.6/+7.3 points over the repair window. Two audits bound the
claim: a single-variable channel ablation shows the retrieval gain does
\emph{not} run through entity-channel fusion votes (metrics bitwise
identical with every entity-dependent channel off), and an LLM-assisted
audit of a seeded 100-link sample puts rebuilt link precision at
\textbf{0.85} (95\% CI [0.77, 0.91])---1.00 on domain terminology, 0.80
on the generic tail that survives repair, a documented boundary rather
than a triumph. We argue the
incident generalizes: wherever an LLM drafts a schema whose aliases then
become the extraction surface, lexicon quality is load-bearing
infrastructure, and it needs its own observability---scale-invariant
link-concentration and entity-name-distribution monitors that standard
pipeline dashboards do not provide by default.
\end{abstract}

\begin{keyword}
knowledge base construction \sep schema quality \sep terminology
extraction \sep entity linking \sep production incident \sep observability
\end{keyword}
\end{frontmatter}

\section{Introduction}
\label{sec:intro}

Large scholarly knowledge platforms increasingly generate their subject
schemas with LLM assistance: a batch tool drafts entity types with a handful
of aliases each, a validator checks structural well-formedness, and the
schema is bound to knowledge bases as configuration (our own batch generator
is the first villain of this story; LLM-for-ontology generation is surveyed
in~\cite{llmonto}). This is attractive at scale---our deployment covers
80{,}855 parsed journal articles (80{,}852 active, 7{,}229{,}225 text
chunks) across two large knowledge bases and several smaller ones---and it
keeps humans out of a loop that would otherwise require domain
lexicographers per discipline.

This paper is about what that design decision does to the entity layer when
the enrichment pipeline extracts mentions by dictionary scan. In our system,
stage-K1 mention extraction is deliberately LLM-free: a schema's
\texttt{entity\_types} main names and aliases constitute the scan
dictionary. The choice is defensible---dictionary scan is cheap enough to
run over 7.2 million chunks, deterministic, idempotent under re-ingest---and
it performed adequately on the knowledge base whose schema had been
hand-evolved over years. It failed catastrophically, and silently, on the
two knowledge bases whose schemas came out of the batch generator.

The failure signature: a 61{,}137-document nuclear-electronics knowledge
base (104112E6) accumulated \textbf{19.2 million entity links onto 575
canonical entities}; the entities were not domain concepts but type
words---``method'', ``parameter'', ``energy''. A 13{,}408-document physics
knowledge base (693372D5) showed the same disease at 2{,}018 entities.
Meanwhile every operational signal was green: enrichment jobs completed
normally, link volume read as growth, and retrieval endpoints answered. The
entity graph had degenerated into a bipartite ``type-to-document'' index
while presenting, to every surface-level check, as a healthy knowledge
graph.

We contribute:

\begin{enumerate}
\item \textbf{A forensic record of entity collapse induced by schema-lexicon
quality}, including the three-layer causal chain
(Section~\ref{sec:collapse}) and an analysis of why standard monitoring is
blind to it---the observable that finally exposes it is entity cardinality
against corpus size, plus the name distribution of top-linked entities.
\item \textbf{A three-source, corpus-attested lexicon generation method}
(Section~\ref{sec:lexicon}) that repairs the affected schemas without manual
lexicography: 2{,}284 transplanted seed terms from a healthy near-domain
library, 7{,}391 corpus-frequency candidates classified by LLM, and bounded
per-type expansion---yielding 3{,}702 real terms for nuclear electronics
(v2.1) and 1{,}279 for physics (v1.1), with 223 legacy cross-type conflicts
repaired in passing.
\item \textbf{A full-chain rebuild at production scale}
(Section~\ref{sec:rebuild}): vertex- and table-scoped teardown across graph
store, relational registry, and search projection; 766{,}501 re-enqueued
jobs; and the pipeline pathologies the rebuild stress test surfaced, each
with a measured fix---most dramatically a 137$\times$ link-throughput
recovery once an $O(N^2)$ alias check, itself amplified by the new richer
lexicon, was cached.
\item \textbf{Before/after evidence and observability recommendations}
(Sections~\ref{sec:after}--\ref{sec:discussion}): entity cardinality
575$\to$1{,}871 with a functioning relation layer (44k typed relations,
1.56M evidence rows); scoped retrieval trajectory +6.6/+7.3 cov@10 with
a channel-level attribution bound (Section~\ref{sec:after}) and an
audited link precision of 0.85 with a 1.00 domain stratum; and four
monitors---led by the scale-invariant top-10 link concentration---that
would have caught the incident on day one.
\end{enumerate}

The larger claim we defend: in dictionary-driven pipelines, \textbf{the
schema lexicon is not descriptive metadata but operative extraction
surface}. Treating it as configuration-file decoration---easily
batch-generated, rarely audited---is a structural mistake that this incident
renders measurable.

\section{Background: dictionary-driven enrichment}
\label{sec:background}

Our enrichment pipeline (four stages per document: mentions $\to$ links
$\to$ relations $\to$ graph) extracts mentions in stage K1 by normalized
dictionary scan: every schema alias match over chunk text yields a
\texttt{raw\_mention} with an entity-type hint; stage K2 links mentions to
canonical entities by normalized form + type; K3 lifts selected mention
pairs into typed relations with an LLM pass; K4 projects the per-document
entity subgraph into the graph store and search-side entity index.
Retrieval consumes all four: an entity channel seeds personalized PageRank,
a co-occurrence channel, and entity-vector similarity all trace back to
linked mentions.

Why no LLM in K1? Cost and determinism. At 7.2M chunks, a per-chunk LLM
extraction pass is a material recurring expense and a re-ingest hazard
(non-determinism breaks idempotent rebuild); a scan is $O(\mathrm{corpus})$
with zero marginal model cost. The design bets everything on one input: the
dictionary. That bet is invisible in the architecture diagram and total in
its consequences---a point Section~\ref{sec:collapse} makes empirical.

The healthy control in our own deployment: a hydrogeology knowledge base
whose schema was hand-evolved to 3{,}742 genuine domain terms (78 documents
$\to$ 1{,}608 canonical entities, well-distributed). The contrast cases:
schemas produced by our batch tool, whose output per entity type is a small
set of \emph{category} aliases (4--8 per type)---a taxonomy of the subject,
not a terminology of it (the distinction between terminological and
taxonomic vocabulary is classical~\cite{iso1087}; ontology-quality metrics
literature).

Neither half of our pipeline is novel in isolation. Dictionary- or
gazetteer-driven entity recognition is one of the oldest NER
traditions~\cite{nerSurvey}, and corpus-statistical term recognition with
linguistic evidence goes back to C-value/NC-value~\cite{cvalue}---our
corpus-candidate classification source (Section~\ref{sec:lexicon}) is that
automatic-term-recognition program reborn with an LLM in the classifier
seat. Likewise, detecting and repairing degraded knowledge-graph content is
an established concern~\cite{paulheim}, and production data-validation
systems routinely gate model inputs with schema- and statistics-level
checks~\cite{tdd}. What is missing from that prior art, and what this
incident supplies, is the specific failure locus: when the dictionary
\emph{is} the extraction model, a structurally valid but terminologically
hollow schema passes every existing gate---term-recognition research
assumes the candidate pool is real corpus vocabulary, and data-validation
systems check schema shape and distribution drift, not whether a lexicon's
entries are domain-specific. The monitors we propose
(Section~\ref{sec:discussion}) close exactly that gap.

\section{The collapse, forensically}
\label{sec:collapse}

\textbf{The observation.} During an unrelated storage investigation
(2026-10-08), a routine entity-quality sample on 104112E6 returned type
words where domain concepts were expected. The full picture
(Table~\ref{tab:collapse}):

\begin{table}[t]
\centering
\caption{The collapse, measured.}
\label{tab:collapse}
\begin{tabular}{lllll}
\toprule
KB & Docs & Entities & Links & Schema provenance \\
\midrule
104112E6 (nuclear electronics) & 61{,}137 & \textbf{575} & \textbf{{\raisebox{1pt}{$\scriptstyle\sim$}}19.2M} & batch-generated \\
693372D5 (physics) & 13{,}408 & \textbf{2{,}018} & $\approx$10$^6$ (est.) & batch-generated \\
hydrogeology (control) & 78 & 1{,}608 & normal & hand-evolved \\
\bottomrule
\end{tabular}
\end{table}

61{,}137 documents carrying ${\sim}$19.2M links is ${\sim}$315 links per
document on average, and ${\sim}$33{,}000 links per entity: every document
``mentions'' the type vocabulary constantly. The graph is a bipartite
type--document index wearing a knowledge graph's clothes.

\textbf{Layer 1 --- mechanism.} K1 scans for schema aliases. When the alias
set for type \texttt{method} contains common words (``method'',
``approach'', and their CJK equivalents), every occurrence of those words in
7.2M chunks becomes a mention, and K2's normalized-form+type linkage folds
them onto one canonical entity per (form, type) pair. The mechanism executed
flawlessly on garbage input.

\textbf{Layer 2 --- provenance.} The dictionaries were the schemas. The
batch generator, prompted to ``draft entity types with aliases for subject
X'', produced exactly that: type-level category vocabulary. Nothing in its
validator checks whether aliases are \emph{domain-specific}---structural
validity is the gate, and category words are structurally valid.

\textbf{Layer 3 --- the control.} Hydrogeology's 3{,}742 hand-evolved terms
over 78 documents yield 1{,}608 entities ${\approx}$20.6 per document of
genuine concepts. Collapse correlates with lexicon nature, not corpus size
or pipeline version.

\textbf{Why nothing fired.} Job-level observability was all green (jobs
DONE, throughput nominal, error rate zero). Link volume---if anything---read
as \emph{success} (19.2M links!). The two observables that expose the
collapse are (a)~entity cardinality per document, which dropped two orders
of magnitude below the healthy control, and (b)~the name distribution of
top-linked entities, which is pure type vocabulary. No standard pipeline
ships either. This blindness generalizes to any dictionary-driven system
whose dashboards count jobs and rows.

\section{Lexicon generation, once and correctly}
\label{sec:lexicon}

The repair requirement: produce genuine domain terminology at scale, without
per-subject lexicographers, without hallucinated terms---and make it a
repeatable tool, not a one-off patch. Our generator implements three
sources, all gated:

\begin{enumerate}
\item \textbf{Seed transplant (2{,}284 terms).} A healthy near-domain
library (the RADAI corpus knowledge base, same academic register) donates
its terms with types. Transplant is the highest-precision source; corpus
attestation below still applies.
\item \textbf{Corpus-candidate classification (7{,}391 candidates).}
High-frequency n-grams from the target corpus's own text are
proposed---concretely, the top 8{,}000 n-grams (unigrams and bigrams) by
chunk frequency over a 3{,}000-chunk-per-KB sample, keeping candidates with
frequency $\geq$5 after a boilerplate and stop-word filter (section
vocabulary such as \emph{acknowledgements/approach/using}; closed-class
English words). An LLM classifies each surviving candidate into a schema
entity type or rejects it. The corpus
floor is the anti-hallucination gate: a term the corpus never attests is
not admitted, whatever the LLM thinks of it.
\item \textbf{Per-type expansion ($\leq$60 per type).} For each entity type
the LLM proposes additional domain terms, given the type description and
the transplanted seeds as in-context anchors; the same corpus-attestation
gate applies. The bound keeps any single type from flooding.
\end{enumerate}

Generation parameters, for reproducibility: the classifier/expander is
Qwen3.6-35B-A3B (AWQ 4-bit) served via vLLM behind a round-robin router,
temperature 0.2, $\leq$4{,}000 output tokens per call, three retries on
empty or non-JSON responses; the entire three-source run for both schemas
is a single idempotent invocation of the generator.

Quality gates beyond attestation: \emph{catch-all types}
(\texttt{term}, \texttt{term\_cn}, \texttt{model\_or\_code},
\texttt{concept}, \texttt{question}) accept no classified noise---they are
the natural sink of classification uncertainty and the main entry route for
taxonomy garbage; a boilerplate filter removes layout/section vocabulary;
cross-type alias conflicts resolve by a specificity ordering (the more
specific type wins); and the schema validator runs as a final structural
gate before any version bump. The same pass repaired 223 legacy cross-type
conflicts in existing schemas.

Output: \textbf{nuclear\_electronics v2.1 = 3{,}702 terms; physics v1.1 =
1{,}279 terms}. 104112E6 rebinds to nuclear\_electronics (its corpus is
nuclear-domain journals). Both versions are idempotent re-runnable---the
tool is now part of the schema lifecycle, not an incident artifact.

An honest note on the asymmetry (physics 1{,}279 versus nuclear 3{,}702):
the physics corpus is smaller and its seed transplant thinner; v1.1 is a
floor, not a ceiling, and the runbook treats small-physics libraries as
follow-up work (Section~\ref{sec:limitations}).

\section{Rebuild at scale}
\label{sec:rebuild}

Repairing live knowledge bases requires tearing down the collapsed entity
layer without stopping ingestion or retrieval. The rebuild driver executes
per knowledge base: delete graph-store vertices with edges; delete the four
relational tables (mentions, links, relations, graph-projection queue rows)
in bounded chunks; delete the search-side entity projection by query; then
re-enqueue all documents through the four stages---\textbf{766{,}501 jobs}
across the two affected libraries. The rebuild ran against the production
fleet (five multi-worker enrich instances, one links-dedicated consumer per
KB) while normal retrieval served.

The rebuild doubled as the system's first full-corpus stress test of the
\emph{rich-lexicon} configuration, and four pathologies surfaced---each now
a fixed, committed, measured lesson:

\begin{itemize}
\item \textbf{B90 --- unbounded projection join.} Graph projection
originally joined each document against \emph{whole-corpus} entities and the
\emph{whole-corpus} link table; at 80k documents and millions of links a
single projection took 25+ minutes, and 80 concurrent consumers saturated a
112-core MySQL (107 cores, load 500). Scoping the join to the document's own
mentions (an index-ordered path) is semantically equivalent---an entity
projects through documents that mention it---and took projection to
seconds, database to 29 cores.
\item \textbf{B91/B93 --- link-insertion deadlocks.} Two deadlock classes in
link insertion: unordered INSERT-SELECT acquiring gap locks in interleaved
order (fixed by pinning insertion order to \texttt{mention\_id}), and
per-action round-trip inserts (fixed by folding all actions into a single
UNION-ALL derived-table INSERT-SELECT). Deadlock error 1213 went to zero.
\item \textbf{B94 --- the $O(N^2)$ the new lexicon amplified.} Alias
canonicalization checked every acronym pair quadratic in alias count---95
of 101 seconds per links task, burning the GIL while 16 threads queued.
Caching canonical forms by schema signature (lexicon version + alias count)
took it to 0.05\,s. Net link throughput: \textbf{8/min $\to$ 1{,}100/min
(137$\times$)}---the collapse-era lexicon was so small it hid this bug for
months; the \emph{repair} exposed it in an afternoon.
\end{itemize}

The last point deserves emphasis in prose, not just as a bullet: the
collapsed lexicon (575 de-facto terms) masked an $O(N^2)$ hot path because
$N$ was tiny. Correcting the lexicon did not merely improve extraction---it
changed the computational regime of the pipeline and surfaced latent
engineering debt. Lexicon quality is load-bearing in the most literal sense.

The rebuild also supplied its own fifth pathology, and it is the one that
most cleanly demonstrates our observability thesis. The B93
single-statement link writer shipped with a parameter-order bug: the SQL's
first placeholder was the knowledge-base id (\texttt{UNHEX(\%s)}), but the
call bound the term form there first---so every row carried
\texttt{UNHEX('alloy')} $=$ NULL as its kb\_id, violating the NOT NULL
column, and \textbf{INSERT IGNORE silently swallowed the entire insert}.
Jobs completed. Entities were created. Throughput read 1{,}100 jobs per
minute. Only one observable was wrong: the link registry stopped
growing---for seven hours. We detected it via row-count reconciliation
against job completions, fixed the binding, reset 51{,}459 $+$ 24{,}372
silently-empty jobs, and measured recovery at +611k rows in the first 60
seconds. Job-level DONE is not row-level truth; any bulk writer guarded by
INSERT IGNORE needs an affected-rows reconciliation, and pipeline
dashboards need a target-table growth rate beside the job queue.

Rebuild final state. All four stages drained across both knowledge bases:
mentions, links, relations, graph projection. The repaired entity layer,
measured: \textbf{KB1 1{,}871 entities carrying 25{,}097{,}704 links (from
575 / ${\sim}$19.2M collapsed), 44{,}043 relations with 1.56M evidence rows
(from effectively zero); KB2 2{,}730 / 5{,}009{,}755 / 12{,}774.}
Spot-checked canonical entities (germanium detector, hpge detector,
neutron, preamplifier, scintillation detector) all project correctly into
the graph store. Retrieval on KB1 answers in 0.6--1.3\,s warm with all
channels active.

\section{Before/after}
\label{sec:after}

\textbf{Entity layer.} The headline repair (all numbers from the production
database, reproducible SQL in the released evidence pack;
Table~\ref{tab:after}):

\begin{table}[t]
\centering
\caption{Before/after, entity layer.}
\label{tab:after}
\small
\begin{tabularx}{\textwidth}{lXXX}
\toprule
Metric & Collapsed & Rebuilt & Control (hydro) \\
\midrule
KB1 documents & 61{,}137 & 61{,}137 & 78 \\
KB1 canonical entities & \textbf{575} & \textbf{1{,}871} (3.25$\times$) & 1{,}608 \\
KB1 entity links & ${\sim}$19.2M (type words) & \textbf{25{,}097{,}704} (terms) & normal \\
KB2 entities / links & 2{,}018 / $\approx$10$^6$ (est.) & \textbf{2{,}730 / 5{,}009{,}755} & --- \\
KB1 typed relations & ${\sim}$0 effective & \textbf{44{,}043} $+$ 1.56M evid.\ rows & --- \\
\bottomrule
\end{tabularx}
\end{table}

The relation layer did not exist in any usable form before the repair; it
now carries schema-typed, LLM-adjudicated relations with per-relation
evidence quotes, and is what the retrieval graph channel traverses.

\textbf{Name quality, honestly.} The top-linked entities after repair are
\texttt{be} (2.26M links), \texttt{alpha}, \texttt{gamma}, \texttt{beta},
\texttt{method}---generic words and short ambiguous symbols still dominate
the link head. The corpus-attestation gate can only guarantee that a term
\emph{occurs}; it cannot distinguish a high-frequency common word from a
domain term. The \texttt{be} case is instructive because it is not even a
ranking problem but a normalization one: the element symbol \emph{Be} and
the English auxiliary verb \emph{be} are distinct in the source text, but
K1's scan normalizes case (ASCII terms match case-insensitively with word
boundaries), so every auxiliary ``be'' in 7.2M chunks links to the
beryllium entity. A case-sensitive allowlist for short chemical symbols
plus an English stop-word exclusion at admission would remove the bulk of
this mass; \texttt{alpha}/\texttt{gamma}/\texttt{beta} are harder---genuine
nuclear-spectrum symbols that also serve as generic Greek-letter
variables---and need context-type constraints, not stop-word lists.
Genuine terminology (hpge detector, neutron, preamplifier, \ldots) is
present and correctly linked and projected---but the repair relocated the
lexicon-quality problem from ``no terminology at all'' to ``terminology plus
a generic tail.'' We consider this the honest headline of the incident:
repairing a collapsed lexicon is one battle; lexicon hygiene is a permanent
operational discipline (Section~\ref{sec:discussion}'s monitors exist
precisely for this).

\textbf{Retrieval trajectory, with a channel-level attribution bound.} On
the scoped evaluation slices (same runner, same corpus scope) the system
moved from the collapsed-era run (2026-09-24) to the rebuilt state
(2026-10-09): associative queries cov@10 0.533$\to$\textbf{0.599} (+6.6pt,
n=48), single-intent queries 0.603$\to$\textbf{0.675} (+7.3pt, n=39) with
MRR $-$2.0/+5.0pt respectively. The two-week window also carried
retrieval-configuration evolution (document-fusion re-ranking, bilingual
alias expansion, kNN search-depth tuning), so the trajectory measures
system evolution \emph{including} the lexicon repair. To bound the
attribution we ran a single-variable channel ablation under the identical
post-rebuild configuration (in-process planner, same corpus scope, same
questions): a \emph{no-entity} arm with every entity-dependent channel
disabled (co-occurrence, graph, entity-vector, PPR) against the full
configuration. Result: \textbf{metrics are bitwise identical on both
families}---cov@10/R@10/MRR unchanged (nea 0.599/0.938/0.735, ner
0.670/0.923/0.770; the only movement anywhere is nea cov@20 $-$1.0pt, 95\%
CI [$-$3.1, 0.0])---while p50 latency falls from 1.46\,s to 0.52\,s. The
co-occurrence channel fired on only 6/48 association queries (its adaptive
gate suppresses it when base retrieval already serves co-mention
documents) and changed no outcome.

The attribution this forces is precise and we state it against interest:
the +6.6/+7.3pt trajectory is \textbf{not} carried by entity channels
casting fusion votes. The entity layer's retrieval-facing value on these
benchmark families runs through the lexical channel's bilingual alias
expansion (entity names as BM25 terms---a mechanism that needs a real
terminology to expand, and which the collapsed lexicon could not feed) and
through structural surfaces (reranker context, knowledge-chain evidence
for generation) documented in the companion system paper; at the fusion
layer itself, entity channels are currently a latency cost with zero
measured recall benefit on these sets. The lexicon repair was necessary
for every entity-consuming surface; the scoped trajectory number alone
cannot be credited to it.

\textbf{Link precision audit (LLM-assisted).} Distributional monitors say
nothing about whether individual links are \emph{correct}. We audited a
seeded, reproducible uniform sample of 100 links from the rebuilt KB1
registry (\texttt{ORDER BY RAND(20261009)}): for each link the mention
span was programmatically verified against the chunk text (span integrity
100/100---offsets are exact), and an LLM judge (Qwen3.6-35B-A3B,
temperature 0) received the $\pm$250-character context with the span
marked and answered a strict two-part rubric: (i)~does the span denote the
same concept as the entity's canonical name (different sense, different
concept, or non-term mention $\Rightarrow$ invalid; debatable $\Rightarrow$
invalid + borderline); (ii)~is the canonical name a domain-specific
technical concept or a generic academic word. Results: \textbf{overall
precision 0.85} (Wilson 95\% CI [0.77, 0.91]); the domain stratum is
\textbf{25/25 $=$ 1.00} (CI [0.87, 1.00]) and the generic stratum
60/75 $=$ \textbf{0.80} (CI [0.70, 0.88]). Every adjudication, span, and
reason string is released. The reading is the name-quality headline
quantified: where the lexicon carries real terminology the dictionary-scan
link layer is essentially perfect, and the residual error mass sits almost
entirely in the generic tail. The audit is LLM-assisted rather than
human-adjudicated---a limitation we keep visible
(Section~\ref{sec:limitations})---but its stratum split is exactly what
the monitor design predicts, and the programmatic span check removes the
cheapest class of audit noise.

\section{Discussion}
\label{sec:discussion}

\textbf{Observability for lexicon quality.} Four cheap monitors would have
caught both incidents at first symptom, in the spirit of production
data-validation systems~\cite{tdd}:
\begin{itemize}
\item[(a)] \emph{Top-10 link concentration (primary, scale-invariant).}
The fraction of all entity links held by the ten most-linked entities. As
a share it is comparable across corpus sizes and disciplines---exactly
what raw cardinality is not. Measured on our deployment: healthy libraries
run at 0.16--0.18 (hydrogeology 0.184, rebuilt physics KB2 0.162); the
rebuilt nuclear library sits at 0.345 with its generic tail; the two
small libraries \emph{before} small-KB treatment ran at 0.398 and 0.768,
and the fully collapsed state is the degenerate limit where a handful of
type words hold ${\sim}$100\% of links. A band alarm (say,
$>$0.35 investigation, $>$0.5 escalation) catches both the catastrophic
collapse and the slow generic-tail drift.
\item[(b)] \emph{Entity cardinality ratio (coarse secondary).} Canonical
entities per ingested document, alarmed against a per-domain band (our
control suggests O(10--20) for scholarly corpora; collapse ran at 0.009).
Cardinality alone is gameable in the other direction---the small-KB
treatment showed category-word aliases \emph{inflating} the count
(Section~\ref{sec:limitations})---which is why it backs up (a) rather
than leading.
\item[(c)] \emph{Top-entity name audit.} The top-K entities by link count
should be domain concepts; a manual or LLM spot check of that list is a
five-minute quarterly control, and the link-precision rubric of
Section~\ref{sec:after} is a ready-made protocol for a deeper sample.
\item[(d)] \emph{Target-table growth reconciliation.} Completed-job
counts reconciled against downstream row growth, mandatory wherever a
bulk writer is guarded by INSERT IGNORE (the silent-swallow incident of
Section~\ref{sec:rebuild} was invisible to job metrics for seven hours
and exposed by exactly this reconciliation).
\end{itemize}
None of the four requires ground truth.

\textbf{LLM extraction versus dictionary scan.} The incident is not an
argument for LLM-everywhere. It is an argument for honoring the input
contract of cheap deterministic components: if extraction is a scan, the
dictionary is the model---and it deserves the same evaluation discipline
(held-out attestation, distributional checks) as a trained model.
Batch-generating that dictionary with an LLM and shipping it on structural
validation alone is the actual failure mode.

\textbf{Schema tooling implications.} The batch schema generator produced
category taxonomy where a terminology was needed; the repair gave the schema
a real lexicon (3{,}702 terms) and folded lexicon generation into the schema
lifecycle as an idempotent tool. What remains open---and what the rebuilt
top-entity distribution makes plain---is that \emph{admission} criteria
(corpus attestation) are necessary but not sufficient: a lexicon also needs
\emph{ranking} criteria that demote high-frequency generic words and short
ambiguous symbols (IDF-style corpus contrast, bilingual alias consistency,
length/type priors). The observation that the collapsed state and the
repaired state differ by a factor of 3.25 in entity cardinality---while both
pass structural validation---is the strongest argument we have for treating
lexicon quality as first-class schema infrastructure with its own acceptance
gates, not as a one-time repair.

\section{Limitations and honesty}
\label{sec:limitations}

\begin{enumerate}
\item \textbf{The generic tail survives the repair}: top-linked entities
remain short generic symbols; admission gates without ranking gates move
the problem, they do not remove it.
\item \textbf{Retrieval attribution is bounded, not eliminated}: the
channel ablation of Section~\ref{sec:after} shows the trajectory does not
run through entity-channel fusion votes, but the full decomposition
(lexicon old/new $\times$ retrieval configuration on/off) would require
re-binding the collapsed schema and re-running the collapsed-era stack---a
re-collapse of a production library we did not authorize for
attribution's sake.
\item \textbf{Coverage}: two knowledge bases rebuilt (61{,}137 $+$
13{,}408 documents) and, after this draft stabilized, the two smaller
diseased libraries (1{,}324 $+$ 119 documents) with the same playbook---where
entity cardinality \emph{contracted} (870$\to$629, 209$\to$153, with
916k/26k links and 4.6k new typed relations): the category-word aliases had
inflated the count, so cardinality is not the quality metric, name
distribution is. The bilingual alias layer showed isolated anomalies (an
entity whose mentions vanished between extraction generations) that we
document but do not fully explain.
\item \textbf{Incident-response cost}: six secondary incidents across the
repair window (one silent INSERT-IGNORE swallow of seven hours, three
half-open-connection wedges, a watchdog that mistook completion for
failure and deleted the freshly rebuilt graph, one missing-index class)
each extended the repair; we report them as findings, not hidden costs.
\item \textbf{No human-adjudicated ground truth} for entity correctness:
the link-precision audit is LLM-assisted (with programmatic span
verification), and entity-quality claims otherwise rest on distributional
and name-sample evidence. A human-adjudicated stratified sample is the
natural next evaluation step.
\end{enumerate}

\section*{Data availability}

The lexicon-generation tool (three-source, corpus-attested), the rebuild
driver, the link-precision audit driver with all 100 adjudications (spans,
verdicts, reasons), the channel-ablation runner, and the observability
queries (top-10 link concentration, cardinality ratio, top-name audit,
row-growth reconciliation) are released at
\url{https://devgit.cdut.edu.cn/chenzhiqian/knowledgenexus} (institutional
GitLab; access on request). Schema lexicons (nuclear\_electronics v2.1,
physics v1.1) ship as versioned assets in the same repository. The journal
corpus is third-party copyrighted; entity-layer aggregate statistics and
name samples sufficient to reproduce the forensic analysis are tabulated
in the released evidence pack.

\section*{Conflict of interest}

The authors declare no competing interests.

\begin{thebibliography}{9}

\bibitem{cvalue} K.~Frantzi, S.~Ananiadou, H.~Mima, Automatic recognition
of multi-word terms: the C-value/NC-value method, International Journal on
Digital Libraries 3~(2) (2000) 115--130. doi:10.1007/s007999900023.

\bibitem{nerSurvey} D.~Nadeau, S.~Sekine, A survey of named entity
recognition and classification, Lingvisticae Investigationes 30~(1)
(2007) 3--26.

\bibitem{paulheim} H.~Paulheim, Knowledge graph refinement: A survey of
approaches and evaluation methods, Semantic Web 8~(3) (2017) 489--508.
doi:10.3233/SW-160218.

\bibitem{tdd} N.~Polyzotis, S.~Roy, S.~E. Whang, M.~Zinkevich, Data
validation for machine learning, in: Proc. MLSys, 2019.

\bibitem{llmonto} A.~S. Lippolis, M.~J. Saeedizade, R.~Keskis{\"a}rkk{\"a},
S.~Zuppiroli, M.~Ceriani, A.~Gangemi, E.~Blomqvist, A.~G. Nuzzolese,
Ontology generation using large language models, arXiv:2503.05388 (2025).

\bibitem{iso1087} ISO 1087: Terminology work and terminology science ---
vocabulary, International Organization for Standardization, 2019.

\end{thebibliography}

\end{document}
